\documentclass[a5paper,10pt]{article}

% https://github.com/InDevRus/filippov-solutions
% Нумерация страниц
\pagenumbering{gobble}

% Настройка полей
\usepackage{geometry}
\geometry{left=1cm}
\geometry{right=1cm}
\geometry{top=1cm}
\geometry{bottom=1.5cm}

% Вставка таблицы
\usepackage{array}
\usepackage{tabularx}

% Инструменты выравнивания формул
\usepackage{amsmath}
\usepackage{mathtools}

% Диакритика в математике
\usepackage{amsfonts}

% Вычисления размеров на ходу
\usepackage{calc}

% Удобные записи производных
\usepackage[italic=true]{derivative}

% Настройки сносок
\usepackage{enotez}

% Длинные знаки векторов
\usepackage{anyfontsize}
\usepackage[b]{esvect}

% Вставка картинок
\usepackage{graphicx}

% Вставка красной строки
\usepackage{indentfirst}
\setlength{\parindent}{1em}

% Условия в командах
\usepackage{xifthen}

% Луа-скрипты
\usepackage{luacode}

% Настройка команд отдельно в математическом и текстовом режимах
\usepackage{mathcommand}

% Для повторения LaTeX кода
\usepackage{multido}

% Конфигурация отступов формул
\delimitershortfall-1sp
\usepackage{mleftright}
\mleftright

% Растягивание объектов
\usepackage{relsize}

% Обтекание текстом
\usepackage{wrapfig}
\usepackage{subcaption}
\usepackage{floatflt}

% Поддержка ввода кириллицы
\usepackage[english, russian]{babel}

% Настройка корневой позиции для компилятора
\usepackage{import}
\providecommand*{\ProjectRootPath}{..}

\import{\ProjectRootPath/preambles/macros/}{theorems}

\import{\ProjectRootPath/preambles/fonts}{font_setup}

% Знак градуса
\usepackage{siunitx}

\import{\ProjectRootPath/preambles/macros/}{operators}
\import{\ProjectRootPath/preambles/macros/}{redefinitions}

\import{\ProjectRootPath/preambles/fonts}{letters_setup}
\import{\ProjectRootPath/preambles/fonts/adjustments}{custom_superscripts}

\import{\ProjectRootPath/preambles/custom}{formatting}
\import{\ProjectRootPath/preambles/custom}{frames}
\import{\ProjectRootPath/preambles/custom}{list_setup}

% TODO: перевести команды на современный стиль объявления

\begin{document}

Для уравнения
$\pow{\lambda}{5} + 2\pow{\lambda}{4} + 4\pow{\lambda}{3} + 6\pow{\lambda}{2} + 5\lambda + 4 = 0$ 
многочлены из критерия Михайлова имеют вид 
$p\br{\xi} = 2\pow{\xi}{2} - 6\xi + 4$ и $q\br{\eta} = \pow{\eta}{2} - 4\eta + 5$.
Корни второго не являются веществнными, так что критерий Михайлова неприменим.

Матрица Гурвица такова:
$H = \begin{pmatrix}
2 & 1 & 0 & 0 & 0
\\ 6 & 4 & 2 & 1 & 0
\\ 4 & 5 & 6 & 4 & 2 
\\ 0 & 0 & 4 & 5 & 6
\\ 0 & 0 & 0 & 0 & 4
\end{pmatrix}$.

$\sub{\Delta}{4} = \det\begin{pmatrix}
2 & 1 & 0 & 0 
\\ 6 & 4 & 2 & 1
\\ 4 & 5 & 6 & 4
\\ 0 & 0 & 4 & 5
\end{pmatrix} 
= 2 \det \begin{pmatrix}
4 & 2 & 1
\\ 5 & 6 & 4
\\ 0 & 4 & 5
\end{pmatrix}
-\det \begin{pmatrix}
6 & 2 & 1
\\ 4 & 6 & 4
\\ 0 & 4 & 5
\end{pmatrix}
= 2 \cdot 26 - 60 = -8$.

По критерию Льенара-Шипара не все $\lambda$ имеют отрицательную вещественную часть.

Предположим, что среди корней присутствует $\sub{\lambda}{m} = ai$, $a \ne 0$. Из предположения следует
$\pow{a}{5} i + 2\pow{a}{4} - 4\pow{a}{3} i - 6\pow{a}{2} i + 5ai + + 4 = 0$,
то есть $\tsys{& \pow{a}{5} - 4\pow{a}{3} + 5a = 0 \\& 2\pow{a}{4} - 6\pow{a}{2} + 4 = 0}$.
Первое из этих уравнений имеет единственное решение $a = 0$, и оно не подходит под второе.

Итак, среди корней нет чисто мнимых, значит по меньшей мере один имеет 
положительную вещественную часть, а нулевое решение неустойчиво.

\end{document}
